\chapter{Valores singulares, traza y clases de Schatten}\label{cap:v3-07}
La norma de operador mide la mayor expansión posible, pero no describe cuántas direcciones contribuyen a ella. Los valores singulares organizan esas expansiones en una sucesión. Su suma y la suma de sus cuadrados producen dos clases fundamentales de operadores: los de clase de traza y los de Hilbert--Schmidt. Probaremos también la completitud de todas las clases de Schatten, desarrollando el ejercicio más extenso de la lectura.

\section{Raíz positiva y descomposición singular}\label{lec:v3-12}
\fuentelectura{12}{Operadores de Hilbert--Schmidt y de clase de traza}
Diremos que $A\in\Bcal(H)$ es positivo si es autoadjunto y $\ip{Ax}x\ge0$ para todo $x$. En particular $T^*T$ es positivo para todo $T$ acotado.

\begin{proposicion}\label{v3:raiz-compacta}
Si $T$ es compacto, existe un único operador positivo $|T|$ tal que $|T|^2=T^*T$. Es compacto y satisface $\norm{|T|x}=\norm{Tx}$.
\end{proposicion}
\begin{proof}
El operador $T^*T$ es compacto positivo. Por el teorema espectral, sus valores propios no nulos, repetidos según multiplicidad, son $a_n>0$, con vectores propios ortonormales $f_n$, y
\[T^*Tx=\sum_n a_n\ip x{f_n}f_n.
\]
La lista puede ser finita; en el caso infinito $a_n\to0$. Definimos
\[|T|x=\sum_n\sqrt{a_n}\ip x{f_n}f_n,
\]
con valor cero en el núcleo. La serie define un operador positivo y las truncaciones convergen en norma porque el supremo de los coeficientes restantes tiende a cero. Es, pues, compacto y su cuadrado es $T^*T$. Además $\norm{|T|x}^2=\ip{T^*Tx}x=\norm{Tx}^2$.

Si $B$ es otra raíz positiva, $B$ conmuta con $T^*T=B^2$. Por ello preserva cada espacio propio $E_a$ para $a>0$, y su restricción positiva a ese espacio finito tiene cuadrado $aI$. El teorema espectral finito obliga a que todos sus valores propios sean $\sqrt a$, luego $B|_{E_a}=\sqrt a I$. En el núcleo de $T^*T$, $\norm{Bx}^2=\ip{B^2x}x=0$, de modo que $B$ se anula. La descomposición ortogonal demuestra $B=|T|$. La unicidad se afirma entre raíces \emph{positivas}; otras raíces pueden existir.
\end{proof}

\begin{definicion}\index{valor singular}
Los valores propios positivos de $|T|$, repetidos según multiplicidad y ordenados de manera no creciente, se llaman \emph{valores singulares} de $T$ y se denotan $s_n(T)$. Si la lista es finita, la completamos con ceros cuando una fórmula lo requiera.
\end{definicion}

\begin{teorema}[Descomposición singular y polar]\label{v3:singulares}\index{descomposición!singular}\index{descomposición!polar}\index{isometría parcial}
Existen familias ortonormales $(f_n),(g_n)$ tales que, con $s_n=s_n(T)>0$,
\begin{align*}
Tx&=\sum_ns_n\ip x{f_n}g_n,& |T|x&=\sum_ns_n\ip x{f_n}f_n,\\
T^*x&=\sum_ns_n\ip x{g_n}f_n,& |T^*|x&=\sum_ns_n\ip x{g_n}g_n.
\end{align*}
Las series de operadores convergen en norma. Además $T=U|T|$, donde $U$ es una isometría de $\cl{\ran|T|}=(\ker T)^\perp$ sobre $\cl{\ran T}$ y vale cero en $\ker T$.
\end{teorema}
\begin{proof}
Escojamos los $f_n$ como base ortonormal de los espacios propios positivos de $|T|$ y definamos $g_n=Tf_n/s_n$. Entonces
\[\ip{g_n}{g_m}=\frac{\ip{T^*Tf_n}{f_m}}{s_ns_m}=\delta_{nm}.
\]
El núcleo de $|T|$ coincide con el de $T$ por la igualdad de normas anterior. Aplicando $T$ a la descomposición de $x$ en ese núcleo y los $f_n$ obtenemos la primera fórmula; la norma de su cola es $\sup_{n>N}s_n$, por ortogonalidad. Para el adjunto, evaluar los apareamientos en la serie da la tercera fórmula. De ella se deduce $TT^*g_n=s_n^2g_n$ y que $TT^*$ se anula en el complemento ortogonal de los $g_n$; la unicidad de la raíz positiva da la cuarta.

La aplicación definida inicialmente por $U(|T|x)=Tx$ está bien definida porque $\ker|T|=\ker T$, y preserva normas. Se extiende por continuidad a la cerradura de esa imagen, con imagen $\cl{\ran T}$, y se define como cero en su complemento ortogonal. Esto prueba la descomposición polar. En general $U$ es una \emph{isometría parcial}, no un unitario como afirma la lectura: el corrimiento compacto ponderado inyectivo en $\ell^2$ tiene parte polar igual al corrimiento unilateral, que no es suprayectivo.
\end{proof}

\section{Hilbert--Schmidt y la independencia de la base}
Una suma de números no negativos sobre un conjunto arbitrario de índices se define como el supremo de sus sumas finitas. Si es finita, sólo un número a lo sumo numerable de términos es no nulo: para cada $m$, sólo puede haber finitos términos mayores que $1/m$.

\begin{proposicion}\label{v3:hilbert-schmidt}\index{operador!de Hilbert--Schmidt}
Para cualquier $T\in\Bcal(H)$, el número
\[\norm T_{\mathrm{HS}}^2=\sum_{e\in\mathcal E}\norm{Te}^2\in[0,\infty]
\]
no depende de la base ortonormal $\mathcal E$. Además $\norm T_{\mathrm{HS}}=\norm{T^*}_{\mathrm{HS}}$. Si es finito, $T$ es compacto y
\[\norm T\le\norm T_{\mathrm{HS}},\qquad
\norm T_{\mathrm{HS}}^2=\sum_ns_n(T)^2.
\]
\end{proposicion}
\begin{proof}
Para bases ortonormales $\mathcal E,\mathcal F$, Parseval y el intercambio de sumas no negativas dan
\[\sum_{e\in\mathcal E}\norm{Te}^2
=\sum_{e\in\mathcal E}\sum_{f\in\mathcal F}|\ip{Te}f|^2
=\sum_{f\in\mathcal F}\sum_{e\in\mathcal E}|\ip e{T^*f}|^2
=\sum_{f\in\mathcal F}\norm{T^*f}^2.
\]
Tomando primero $\mathcal F=\mathcal E$ obtenemos igualdad para $T$ y $T^*$ en una base; combinándola con la identidad para dos bases obtenemos independencia. Parseval proviene de la expansión ortogonal de los vectores: en particular $\ip xy=\sum_{e\in\mathcal E}\ip xe\ip ey$, con convergencia absoluta por Cauchy--Schwarz.

Para $x=\sum_e\ip xe\,e$ y $y$ unitario, Cauchy--Schwarz da
\[|\ip{Tx}y|\le\Bigl(\sum_e|\ip xe|^2\Bigr)^{1/2}
\Bigl(\sum_e|\ip{Te}y|^2\Bigr)^{1/2}
\le\norm x\norm T_{\mathrm{HS}}.
\]
Así $\norm T\le\norm T_{\mathrm{HS}}$. Si esta última norma es finita, para cada $\eps>0$ existe un conjunto finito $F\subset\mathcal E$ con $\sum_{e\notin F}\norm{Te}^2<\eps^2$. La misma estimación aplicada a $T(I-P_F)$ da $\norm{T-TP_F}<\eps$. Como $TP_F$ tiene rango finito, $T$ es compacto. Finalmente, extendemos los vectores singulares $f_n$ a una base con vectores del núcleo; en ella $\norm{Tf_n}^2=s_n^2$ y los demás términos son cero, dando la fórmula de valores singulares.
\end{proof}

\section{Una fórmula variacional para las clases de Schatten}
Para $1\le p<\infty$, definimos
\[\mathcal S_p(H)=\{T\in\Kcal(H):\sum_ns_n(T)^p<\infty\},\qquad
\norm T_p=\Bigl(\sum_ns_n(T)^p\Bigr)^{1/p}.
\]
Escribimos $\mathcal S_2=\mathrm{HS}(H)$ y $\mathcal S_1=\mathrm{TR}(H)$; la norma de esta última también se denota $\norm T_{\mathrm{tr}}$. La siguiente prueba establece de una vez la desigualdad triangular y la completitud, que no se deducen simplemente de la desigualdad triangular para sucesiones, porque los valores singulares de una suma no se suman término a término.

\begin{lema}[Fórmula variacional]\label{v3:schatten-variacional}
Sea $T$ compacto, $1\le p<\infty$ y $q$ el exponente conjugado, con $q=\infty$ si $p=1$. Entonces, admitiendo el valor $+\infty$,
\[\norm T_p=
\sup\left\{\left|\sum_{i=1}^m b_i\ip{Te_i}{h_i}\right|:
\begin{array}{l}(e_i)_{i=1}^m,(h_i)_{i=1}^m\text{ ortonormales},\\
b_i\ge0,\ \norm b_{\ell^q}\le1,\ m<\infty
\end{array}\right\}.
\]
\end{lema}
\begin{proof}
Use la descomposición $Tx=\sum_js_j\ip x{f_j}g_j$. Para las familias finitas del enunciado pongamos
$a_{ij}=|\ip{e_i}{f_j}|\,|\ip{g_j}{h_i}|$. Por Cauchy--Schwarz y Bessel,
\[\sum_ja_{ij}\le1\quad\text{para cada }i,\qquad
\sum_ia_{ij}\le1\quad\text{para cada }j.
\]
Entonces el módulo de la suma está acotado por $\sum_{i,j}b_is_ja_{ij}$. Si $1<p<\infty$, Hölder para esta suma con pesos no negativos da
\[\sum_{i,j}b_is_ja_{ij}
\le\Bigl(\sum_{i,j}b_i^qa_{ij}\Bigr)^{1/q}
\Bigl(\sum_{i,j}s_j^pa_{ij}\Bigr)^{1/p}
\le\norm b_q\norm s_p\le\norm T_p.
\]
Para $p=1$ se usa directamente $b_i\le1$ y $\sum_i a_{ij}\le1$, obteniendo la misma cota. Las sumas infinitas se justifican por límites monótonos de sumas finitas.

Para la cota opuesta, escoja $e_i=f_i$, $h_i=g_i$ para los primeros $m$ valores singulares positivos. Si $p=1$, tome $b_i=1$; obtiene $\sum_{i=1}^ms_i$. Si $p>1$, tome
$b_i=s_i^{p-1}/(\sum_{j=1}^ms_j^p)^{1/q}$, cuando el denominador no sea cero. Entonces $\norm b_q=1$ y la suma es $(\sum_{i=1}^ms_i^p)^{1/p}$. Al aumentar $m$ se alcanza el valor requerido. Si $T=0$, ambas partes son cero.
\end{proof}

\begin{teorema}\label{v3:schatten-banach}\index{Schatten, clases de}
Para cada $1\le p<\infty$, $\mathcal S_p(H)$ es un espacio de Banach con la norma $\norm\cdot_p$, y $\norm T\le\norm T_p$. En particular $\norm\cdot_{\mathrm{HS}}$ y $\norm\cdot_{\mathrm{tr}}$ son normas.
\end{teorema}
\begin{proof}
Los compactos forman un espacio vectorial. La fórmula variacional expresa $\norm T_p$ como un supremo de módulos de funcionales lineales en $T$; por ello es homogénea y subaditiva. Si $S,T\in\mathcal S_p$, esa subaditividad hace finita la norma de $S+T$, así que la clase es un espacio vectorial. El caso $m=1$, $b_1=1$ muestra $\norm T\le\norm T_p$, y de aquí se deduce que sólo el operador cero tiene norma cero.

Sea $(T_n)$ de Cauchy en $\norm\cdot_p$. También lo es en norma de operador, luego converge a un operador $T$; como los compactos son cerrados en esa norma, $T$ es compacto. Para cada funcional finita $F$ del lema, $F(T_n-T_m)\to F(T_n-T)$ cuando $m\to\infty$. Dado $\eps>0$, si $n,m\ge N$ tenemos $|F(T_n-T_m)|\le\eps$, por tanto $|F(T_n-T)|\le\eps$ para $n\ge N$. Tomando el supremo en todas las elecciones finitas, $\norm{T_n-T}_p\le\eps$. En particular $T-T_N\in\mathcal S_p$, y entonces $T\in\mathcal S_p$. Esto prueba la completitud sin suponer de antemano continuidad de cada valor singular.
\end{proof}

\section{Traza y productos de operadores}
Para $u,v\in H$ escribiremos $u\otimes v$ para el operador de rango uno $x\mapsto\ip xv u$. Su único valor singular no nulo es $\norm u\norm v$ si ambos vectores son no nulos: el cuadrado del módulo es el operador $\norm u^2(v\otimes v)$. Por tanto $\norm{u\otimes v}_p=\norm u\norm v$ para todo $p$.

\begin{proposicion}\label{v3:traza-variacional}
Para un compacto $T$,
\[\norm T_{\mathrm{tr}}=
\sup_{\mathcal E,\mathcal F}\sum_{i\in I}|\ip{Te_i}{h_i}|,
\]
donde se consideran bases ortonormales $\mathcal E=(e_i)_{i\in I}$, $\mathcal F=(h_i)_{i\in I}$ con un mismo conjunto de índices. Si $T$ es de clase de traza, la serie
\[\tr T=\sum_{e\in\mathcal E}\ip{Te}e
\]
converge absolutamente y es independiente de la base. Para una descomposición singular,
\[\tr T=\sum_ns_n\ip{g_n}{f_n},\qquad |\tr T|\le\norm T_{\mathrm{tr}}.
\]
\end{proposicion}
\begin{proof}
La cota superior para la suma de módulos sigue del lema~\ref{v3:schatten-variacional} con $p=1$, absorbiendo las fases en los $h_i$ y tomando supremos sobre subconjuntos finitos. Para la cota inferior, fijado $N$ extendemos las familias \emph{finitas} $f_1,\ldots,f_N$ y $g_1,\ldots,g_N$ a bases completas con índices coincidentes en esos primeros lugares; sus complementos tienen la misma dimensión porque se quitó la misma cantidad finita de vectores. La suma es al menos $\sum_{n=1}^Ns_n$. Hacer $N\to\infty$ da la igualdad. Este argumento no necesita extender simultáneamente dos familias infinitas cuyos complementos podrían tener dimensiones distintas.

Para una base $\mathcal E$ y $T$ de traza, la descomposición singular da
\[\sum_e\ip{Te}e=\sum_e\sum_ns_n\ip e{f_n}\ip{g_n}e.
\]
La suma de módulos de los términos dobles es a lo sumo
\[\sum_ns_n\Bigl(\sum_e|\ip e{f_n}|^2\Bigr)^{1/2}
\Bigl(\sum_e|\ip{g_n}e|^2\Bigr)^{1/2}=\sum_ns_n<\infty.
\]
Podemos, pues, intercambiar las sumas. Parseval da $\sum_e\ip e{f_n}\ip{g_n}e=\ip{g_n}{f_n}$, y obtenemos una expresión independiente de $\mathcal E$. La cota de la traza sigue tomando módulos. La linealidad de la traza se obtiene sumando en una misma base las series absolutamente convergentes.
\end{proof}
Si $T$ es además normal, su descomposición espectral y una base de vectores propios dan
\[\tr T=\sum_{\lambda\in\sigma(T)\setminus\{0\}}\lambda\dim E_\lambda,
\]
con convergencia absoluta, pues los valores singulares son los módulos de esos valores propios con sus multiplicidades.

\begin{observacion}[Una igualdad de imágenes que debe corregirse]
La lectura afirma $\ran\sqrt{|T|}=\ran|T|$. En general sólo coinciden sus \emph{cerraduras}, ambas iguales a $(\ker T)^\perp$. Por ejemplo, en $\ell^2$, sea $|T|e_n=n^{-2}e_n$. El vector $y_n=n^{-2}$ pertenece a $\ran\sqrt{|T|}$ porque su preimagen es $(1/n)\in\ell^2$, pero no a $\ran|T|$, pues la única preimagen formal sería $(1,1,\ldots)$. La demostración de independencia de la traza dada arriba evita esta afirmación falsa. Sí es válida la factorización $T=(U\sqrt{|T|})\sqrt{|T|}$, con ambos factores de Hilbert--Schmidt cuando $T$ es de clase de traza, porque los cuadrados de sus valores singulares suman $\sum_ns_n(T)$.
\end{observacion}

\begin{proposicion}\label{v3:ideales-hs-traza}
$\mathcal S_2$ y $\mathcal S_1$ son ideales bilaterales de $\Bcal(H)$ y satisfacen, para $A$ acotado,
\[\norm{AT}_{\mathrm{HS}},\norm{TA}_{\mathrm{HS}}\le\norm A\norm T_{\mathrm{HS}},
\qquad
\norm{AT}_{\mathrm{tr}},\norm{TA}_{\mathrm{tr}}\le\norm A\norm T_{\mathrm{tr}}.
\]
Si $S,T$ son Hilbert--Schmidt, $ST$ es de traza y $\norm{ST}_{\mathrm{tr}}\le\norm S_{\mathrm{HS}}\norm T_{\mathrm{HS}}$.
\end{proposicion}
\begin{proof}
Para Hilbert--Schmidt, en cualquier base,
$\sum_e\norm{ATe}^2\le\norm A^2\sum_e\norm{Te}^2$. La cota para $TA$ se obtiene aplicando la anterior a $A^*T^*$ y usando la igualdad de normas Hilbert--Schmidt de un operador y su adjunto.

Para traza, la descomposición singular $T=\sum_ns_n(g_n\otimes f_n)$ converge en norma de traza: la norma de la cola es $\sum_{n>N}s_n$. Como $\norm{(Ag_n)\otimes f_n}_{\mathrm{tr}}\le\norm A$, la serie $\sum_ns_n(Ag_n\otimes f_n)$ converge en el Banach $\mathcal S_1$ y coincide en norma de operador con $AT$. Esto da la primera cota. De igual modo $TA=\sum_ns_n(g_n\otimes A^*f_n)$ da la segunda.

Finalmente, $ST$ es compacto. Para familias ortonormales finitas $(e_i),(h_i)$,
\[\sum_i|\ip{STe_i}{h_i}|=\sum_i|\ip{Te_i}{S^*h_i}|
\le\Bigl(\sum_i\norm{Te_i}^2\Bigr)^{1/2}
\Bigl(\sum_i\norm{S^*h_i}^2\Bigr)^{1/2}
\le\norm T_{\mathrm{HS}}\norm S_{\mathrm{HS}}.
\]
La fórmula variacional para $p=1$ implica que $ST$ es de traza y da la cota.
\end{proof}

\section{Ejercicios propuestos}
\begin{ejercicio}\label{ej:v3-l12-1}
Demuestre que $\norm\cdot_{\mathrm{HS}}$ y $\norm\cdot_{\mathrm{tr}}$ son normas en sus respectivas clases.
\end{ejercicio}


\begin{ejercicio}\label{ej:v3-l12-2}
Todo operador de Hilbert--Schmidt es compacto.
\end{ejercicio}


\begin{ejercicio}\label{ej:v3-l12-3}
Pruebe $\norm T\le\norm T_{\mathrm{HS}}$.
\end{ejercicio}


\begin{ejercicio}\label{ej:v3-l12-4}
Demuestre que las clases de Hilbert--Schmidt y de traza son ideales bilaterales en $\Bcal(H)$.
\end{ejercicio}


\begin{ejercicio}\label{ej:v3-l12-5}
Pruebe: (i) si $T$ es de traza, $T^*$ también y $\tr(T^*)=\overline{\tr T}$; (ii) para $S,T$ Hilbert--Schmidt, $\tr(ST)=\tr(TS)$; (iii) $\langle S,T\rangle_{\mathrm{HS}}=\tr(ST^*)$ es un producto interno que hace Hilbert a $\mathrm{HS}(H)$.
\end{ejercicio}


\begin{ejercicio}[Todas las clases de Schatten]\label{ej:v3-l12-6}
Para $1\le p<\infty$, pruebe que $\mathcal S_p(H)$ es un espacio vectorial complejo, que $\norm\cdot_p$ es una norma y que ese espacio es completo.
\end{ejercicio}


\section{Ampliación: núcleos integrales y prueba de Schur}\index{Schur, prueba de}
El siguiente criterio desarrolla la nota de Egor Maximenko \emph{La prueba de Schur simplificada para operadores integrales} \cite{maxschur}. Se distinguen los espacios de entrada y salida para evitar invertirlos en la fórmula.

\begin{teorema}[Schur, versión no ponderada]\label{v3:schur}
Sean $(X,\mu),(Y,\nu)$ espacios de medida $\sigma$-finitos y $K:X\times Y\to\C$ medible. Supongamos
\[\esssup_{x\in X}\int_Y|K(x,y)|\,d\nu(y)\le C_1<\infty,
\quad
\esssup_{y\in Y}\int_X|K(x,y)|\,d\mu(x)\le C_2<\infty.
\]
Entonces $Af(x)=\int_YK(x,y)f(y)\,d\nu(y)$ define un operador $L^2(Y,\nu)\to L^2(X,\mu)$ y $\norm A\le\sqrt{C_1C_2}$.
\end{teorema}
\begin{proof}
Tonelli da
\[\int_X\int_Y|K(x,y)||f(y)|^2\,d\nu(y)\,d\mu(x)
\le C_2\norm f_2^2.
\]
Así la integral interior es finita para casi todo $x$. Cauchy--Schwarz aplicado a $|K|^{1/2}$ y $|K|^{1/2}|f|$ implica allí integrabilidad absoluta de $Kf$ y
\[|Af(x)|^2\le\left(\int_Y|K(x,y)|\,d\nu(y)\right)
\left(\int_Y|K(x,y)||f(y)|^2\,d\nu(y)\right).
\]
Integramos y usamos las dos cotas en su orden:
\[\norm{Af}_2^2\le C_1\int_Y|f(y)|^2\left(\int_X|K(x,y)|\,d\mu(x)\right)d\nu(y)
\le C_1C_2\norm f_2^2.
\]
Esto prueba existencia casi en todas partes, independencia del representante y acotación. No debe introducirse un segundo factor $C_1C_2$ al reescribir la integral intermedia, como ocurre en una línea de la nota complementaria.
\end{proof}

\begin{proposicion}\label{v3:nucleo-hs}
Si $K\in L^2([0,1]^2)$, el operador $A_Kf(x)=\int_0^1K(x,y)f(y)\,dy$ es Hilbert--Schmidt y
\[\norm{A_K}_{\mathrm{HS}}=\norm K_{L^2([0,1]^2)}.
\]
\end{proposicion}
\begin{proof}
Por Cauchy--Schwarz en $y$ y después integración en $x$, $\norm{A_Kf}_2\le\norm K_2\norm f_2$. Sea $(e_n)$ una base ortonormal de $L^2[0,1]$. Para casi todo $x$, $K(x,\cdot)\in L^2$ por Fubini y
\[\sum_n\left|\int_0^1K(x,y)e_n(y)\,dy\right|^2
=\norm{K(x,\cdot)}_2^2,
\]
porque $(\cl{e_n})$ también es una base ortonormal y se aplica Parseval. Tonelli permite integrar esta suma no negativa:
\[\sum_n\norm{A_Ke_n}_2^2=\int_0^1\norm{K(x,\cdot)}_2^2\,dx=\norm K_2^2.
\]
La proposición~\ref{v3:hilbert-schmidt} da la conclusión y también la compacidad.
\end{proof}

\begin{ejercicio}[Complemento: núcleo acotado en medida finita]\label{ej:v3-ampl-schur}
Si $\mu(X),\nu(Y)<\infty$ y $|K(x,y)|\le M$ casi en todas partes, pruebe $\norm A\le M\sqrt{\mu(X)\nu(Y)}$.
\end{ejercicio}


\begin{ejercicio}[Complemento: tamaño Hilbert--Schmidt de Volterra]\label{ej:v3-ampl-volterra-hs}
Calcule la norma Hilbert--Schmidt de $V$ en $L^2[0,1]$ y determine si $V^*V$ es de traza.
\end{ejercicio}


Las clases de Schatten proporcionan normas que distinguen diferentes grados de compacidad. La traza recupera una suma de valores propios para compactos normales, mientras que la clase Hilbert--Schmidt convierte el espacio de operadores en un Hilbert. Los criterios integrales conectan estas construcciones con ejemplos calculables; más adelante, las topologías débiles permitirán estudiar límites que no convergen en ninguna de estas normas.
